\documentclass[11pt]{article}
\usepackage[margin=1.1in]{geometry}
\usepackage[T1]{fontenc}
\usepackage{microtype}
\usepackage{amsmath,amssymb,amsthm}
\usepackage{booktabs}
\usepackage{tikz}
\usetikzlibrary{trees}
\usepackage{xurl}
\usepackage[hidelinks]{hyperref}
\hypersetup{pdftitle={Rigidity of Pattern-Avoiding Breadth-First Reading Words of Increasing Trees},pdfauthor={Victor Bona},pdfsubject={Enumerative combinatorics; pattern avoidance in ordered trees},pdfkeywords={pattern avoidance, increasing trees, k-ary heaps, breadth-first search, OEIS, Lean}}
\newtheorem{theorem}{Theorem}[section]
\newtheorem{lemma}[theorem]{Lemma}
\newtheorem{corollary}[theorem]{Corollary}
\newtheorem{proposition}[theorem]{Proposition}
\theoremstyle{definition}
\newtheorem{definition}[theorem]{Definition}
\newtheorem{remark}[theorem]{Remark}
\newtheorem{question}[theorem]{Question}
\newcommand{\Av}{\operatorname{Av}}
\newcommand{\UB}{\mathcal U}
\newcommand{\FB}{\mathcal F}
\newcommand{\HP}{\mathcal H}
\newcommand{\bfs}{\operatorname{bfs}}
\newcommand{\seq}[1]{\href{https://oeis.org/#1}{\texttt{#1}}}
\title{Rigidity of Pattern-Avoiding Breadth-First Reading Words\\ of Increasing Trees}
\author{Victor Bona\\[2pt]\small Independent Researcher, Blumenau, Brazil\\\small\texttt{victor.bona@hotmail.com}}
\date{September 2026}
\begin{document}
\maketitle

\begin{abstract}
We study permutations obtained by reading increasing ordered trees in breadth-first order. For every integer $k\ge2$, we prove two structural results. A $312$-avoiding permutation is realizable on a tree of maximum outdegree $k$ if and only if it is realizable on the complete $k$-ary heap shape. A $231$-avoiding permutation of length congruent to $1$ modulo $k$ is realizable on a tree of maximum outdegree $k$ if and only if it is realizable on a full $k$-ary tree. Both proofs use the nondecreasing sequence of parent positions in breadth-first order. The first follows from a capacity inequality; the second groups nonroot positions into blocks of $k$ and preserves their labels by a pattern-forced inequality. For binary trees these results prove three identities between OEIS sequences, including the conjectured equality of \seq{A245899} and \seq{A246747}. The analogous heap collapse for $321$ fails first at length $4$, whereas full binary realizability first diverges at odd length $11$, with $8095$ unary-binary words and $8048$ full binary words. We provide $28$ additional sequence entries relative to the recorded baseline, computation programs, and the complete set of $47$ counterexamples. The binary tree-level results and the general parent-sequence arguments are verified in Lean~4. We also deduce exponential growth rate $4$ for all three patterns from Defant's heap-growth theorem.
\end{abstract}

\section{Introduction}\label{sec:intro}

Pattern avoidance in permutations, originating in the study of stack sorting, provides a framework for counting combinatorial objects subject to relative-order restrictions \cite{Knuth1997,Bona2012,Kitaev2011}. A reading map transfers these restrictions to labeled structures: one orders their vertices, reads the labels, and requires the resulting permutation to avoid a fixed pattern. Levin, Pudwell, Riehl, and Sandberg studied this question for complete $k$-ary heaps \cite{LPRS2016}. Related work considers linear extensions of posets \cite{AERRSV2016,Defant2019} and increasing unary-binary trees whose breadth-first words avoid $213$ \cite{BFMR2017}.

This paper counts distinct reading words across a family of tree shapes. An increasing ordered tree has distinct labels $1,\ldots,n$, with every child larger than its parent, and the children of each vertex are ordered. Its breadth-first word lists the labels level by level, from left to right. Different labeled trees can produce the same word. Consequently, counting these words is a different problem from counting trees, and an equality of word sets is stronger than an equality of their cardinalities.

For a pattern $\tau$, let $\UB_n^{(k)}(\tau)$ denote the set of $\tau$-avoiding words realizable on increasing ordered trees of maximum outdegree $k$. Let $\FB_n^{(k)}(\tau)$ denote the corresponding set for full $k$-ary trees, in which each outdegree is $0$ or $k$, and let $\HP_n^{(k)}(\tau)$ denote the words on the complete $k$-ary heap shape. Our main results are
\begin{align}
\UB_n^{(k)}(312)&=\HP_n^{(k)}(312),\label{eq:introA}\\
\UB_n^{(k)}(231)&=\FB_n^{(k)}(231)\qquad(n\equiv1\pmod{k}).\label{eq:introB}
\end{align}
The first equality replaces a union over admissible shapes with one prescribed shape. The second replaces a degree bound with an exact degree condition at every internal vertex. Neither asserts uniqueness of an original realization.

The common proof device is the parent sequence. Number vertices by their positions in the breadth-first traversal. The parents of positions $2,\ldots,n$ form a nondecreasing sequence in which each index occurs at most $k$ times. This representation gives a short capacity argument for \eqref{eq:introA} and an explicit block construction for \eqref{eq:introB}. Both transformations preserve the word itself. Given a realization, the second construction takes linear time.

The binary specializations concern six word-counting sequences recorded in the OEIS: \seq{A245898}, \seq{A245899}, and \seq{A245900} for unary-binary trees and patterns $231$, $312$, and $321$, and \seq{A245901}, \seq{A245902}, and \seq{A245903} for full binary trees \cite{OEIS}. The entry \seq{A246747} counts complete binary heaps avoiding $231$ and, equivalently, $312$; its recorded tentative equality with \seq{A245899} follows from \eqref{eq:introA}. The two full binary identities follow from \eqref{eq:introA} and \eqref{eq:introB}. We distinguish these proofs from their historical annotations: the OEIS records use tentative language for the $312$ relationships, whereas the $231$ odd-index relationship is stated affirmatively without a proof on the entry page.

The $321$ pattern separates the two conclusions. The heap condition already fails for the realizable word $1423$. Full binary realizability agrees with unary-binary realizability at odd lengths through $9$, but fails at length $11$. We give a direct obstruction for one word and an exhaustive computation of all $47$ missing words. We also extend the recorded sequence data and determine the exponential growth rate by applying the general theorem of Defant \cite{Defant2019}. A pattern-free promotion lemma, useful in the original binary proof, is retained as a separate exchange principle in Appendix~\ref{sec:cascade}.

\section{Words, trees, and parent sequences}\label{sec:prelim}

Write $[n]=\{1,\ldots,n\}$, where $n\ge1$, and let $w=w_1\cdots w_n$ be a permutation of $[n]$. An occurrence of $231$ consists of positions $a<b<c$ with $w_c<w_a<w_b$; an occurrence of $312$ has $w_b<w_c<w_a$; an occurrence of $321$ has $w_c<w_b<w_a$. A word avoids a pattern when it has no such occurrence. We use $\Av_n(\tau)$ for the corresponding avoidance class.

A rooted ordered tree is increasing if the labels increase along every parent-child edge. Its root then has label $1$. Throughout, $k\ge2$ is an integer. Maximum outdegree $k$ means that every vertex has at most $k$ children; full $k$-ary means that every vertex has either zero or exactly $k$ children. A full $k$-ary tree with $m$ internal vertices has $km$ edges and hence $km+1$ vertices. Its size is therefore congruent to $1$ modulo $k$.

The complete $k$-ary shape fills each level from left to right before starting the next. In one-based breadth-first indexing its parent map is
\begin{equation}
q_k(i)=\left\lceil\frac{i-1}{k}\right\rceil,\qquad 2\le i\le n.\label{eq:heap}
\end{equation}
A permutation labels this shape increasingly exactly when $w_{q_k(i)}<w_i$ for every nonroot position. For $k=2$, $q_2(i)=\lfloor i/2\rfloor$. The complete shape is full $k$-ary when $n\equiv1\pmod{k}$.

\begin{definition}
A parent sequence on $[n]$ is a sequence $(p_2,\ldots,p_n)$ satisfying $1\le p_i<i$ and $p_i\le p_{i+1}$. Its capacity is at most $k$ if each index occurs at most $k$ times. It is full of order $k$ if every occurring index occurs exactly $k$ times. It is compatible with $w$ if $w_{p_i}<w_i$ for all $i\ge2$.
\end{definition}

\begin{proposition}[Parent-sequence representation]\label{prop:parents}
A permutation $w$ is a breadth-first word of an increasing ordered tree of maximum outdegree $k$ if and only if it has a compatible parent sequence of capacity at most $k$. The same equivalence holds between full $k$-ary trees and compatible full parent sequences of order $k$.
\end{proposition}
\begin{proof}
In a breadth-first traversal, the children of each parent form a consecutive block, and the blocks occur in parent order. Thus parent indices are nondecreasing. Each parent precedes its children, its multiplicity is its outdegree, and an increasing labeling gives compatibility.

Conversely, attach vertex $i$ to vertex $p_i$, ordering children by their indices. Every parent index is smaller than its child index, so repeatedly following parents reaches vertex $1$. These edges define a rooted ordered tree. To check its breadth-first order, start a queue with vertex $1$ and append the children of each removed vertex. Nondecreasing parent indices ensure that the concatenation of these child blocks is exactly $2,\ldots,n$. More explicitly, when vertex $j$ is next, every vertex of index at most $j$ has already been enqueued, since its parent has smaller index; the pending indices and subsequent child blocks remain increasing. Induction therefore gives the traversal $1,\ldots,n$. Multiplicities give the required degree condition, and compatibility gives the increasing labeling.
\end{proof}

The distinction between words and trees is already visible at five vertices. The word $12435$ has the unary-binary realization with levels $(1),(2,4),(3,5)$ and edges $2\to3$, $4\to5$ between the last two levels. It also has a full binary realization with the same levels in which both $3$ and $5$ are children of $2$. The word is counted once in either word class.

\section{Heap collapse for 312}\label{sec:A}

\begin{lemma}[Capacity bound]\label{lem:capacity}
For a parent sequence of capacity at most $k$, one has $q_k(i)\le p_i$ for every $i\ge2$.
\end{lemma}
\begin{proof}
By monotonicity, the $i-1$ entries $p_2,\ldots,p_i$ all belong to $[p_i]$. Each possible value occurs at most $k$ times, so $i-1\le kp_i$. Taking a ceiling gives the result.
\end{proof}

\begin{theorem}[Heap collapse]\label{thm:A}
For every $k\ge2$ and $n\ge1$,
\[
\UB_n^{(k)}(312)=\HP_n^{(k)}(312).
\]
Equivalently, a $312$-avoiding permutation is realizable with maximum outdegree $k$ if and only if $w_{q_k(i)}<w_i$ for every $2\le i\le n$.
\end{theorem}
\begin{proof}
The complete heap shape has maximum outdegree $k$, so the reverse inclusion is immediate. For the forward inclusion, take a compatible parent sequence and fix $i\ge2$. Write $q=q_k(i)$ and $p=p_i$. Lemma~\ref{lem:capacity} gives $q\le p<i$, and compatibility gives $w_p<w_i$. If $q=p$ the desired inequality follows. If $q<p$ and $w_q>w_i$, the positions $q<p<i$ have values $w_q>w_i>w_p$, forming $312$. This is impossible. Distinctness of the labels therefore gives $w_q<w_i$.
\end{proof}

The proof supplies the complete shape directly. Once avoidance of $312$ is known, realizability is decided by the $n-1$ comparisons in \eqref{eq:heap}; constructing the canonical realization takes $O(n)$ time. This statement does not include the cost of checking pattern avoidance.

\begin{corollary}\label{cor:full312}
If $n\equiv1\pmod{k}$, then
\[
\UB_n^{(k)}(312)=\FB_n^{(k)}(312)=\HP_n^{(k)}(312).
\]
\end{corollary}
\begin{proof}
At these sizes the complete shape is full $k$-ary. Combine this fact with Theorem~\ref{thm:A} and the inclusion $\FB_n^{(k)}(312)\subseteq\UB_n^{(k)}(312)$.
\end{proof}

\section{Full-degree collapse for 231}\label{sec:B}

The relevant consequence of $231$ avoidance is forward compatibility: if positions $p<j$ satisfy $w_p<w_j$, then every later value $w_i$, $i>j$, also exceeds $w_p$. Otherwise $w_i<w_p<w_j$ at positions $p<j<i$ would form $231$.

\begin{theorem}[Full-degree collapse]\label{thm:B}
For every $k\ge2$ and $n\equiv1\pmod{k}$,
\[
\UB_n^{(k)}(231)=\FB_n^{(k)}(231).
\]
Given a realization in the left-hand class, a realization in the right-hand class can be constructed in $O(n)$ time.
\end{theorem}
\begin{proof}
Only the forward inclusion needs proof. The singleton word is immediate, so suppose $n=km+1$ with $m\ge1$. Take a compatible parent sequence $(p_2,\ldots,p_n)$ of capacity at most $k$. Partition the nonroot positions into the blocks
\[
B_t=\{2+tk,\ldots,1+(t+1)k\},\qquad 0\le t<m,
\]
and put $j_t=2+tk$. Assign every position in $B_t$ to the original parent $p_{j_t}$ of its first position. Thus the new parent sequence is
\begin{equation}
p'_i=p_{j_t}\qquad(i\in B_t).\label{eq:regroup}
\end{equation}

The selected parents strictly increase. If $p_{j_t}=p_{j_{t+1}}$, monotonicity would make all $k+1$ entries from $p_{j_t}$ through $p_{j_{t+1}}$ equal, violating capacity. Consequently $(p'_i)$ is nondecreasing, and every occurring parent appears exactly $k$ times. Also $p'_i=p_{j_t}<j_t\le i$.

Compatibility holds at $i=j_t$ by the original edge. For $i>j_t$ in the same block, a failure would give $w_i<w_{p_{j_t}}<w_{j_t}$ at positions $p_{j_t}<j_t<i$, an occurrence of $231$. Thus $w_{p'_i}<w_i$ throughout. Proposition~\ref{prop:parents} reconstructs an increasing full $k$-ary tree with word $w$.

Reading the original parent array, evaluating \eqref{eq:regroup}, and building the new child lists each take linear time.
\end{proof}

For example, the $231$-avoiding word $1243567$ has a realization of maximum outdegree $3$ with parent sequence $(1,1,2,3,3,4)$. Grouping the six nonroot positions into two blocks gives $(1,1,1,3,3,3)$. The resulting full ternary tree has levels $(1)$, $(2,4,3)$, and $(5,6,7)$, with the final three labels all children of $4$.

For $k=2$ this is a parity collapse. Unlike Theorem~\ref{thm:A}, it does not prescribe one shape for all words. At length $4$, for example, the $231$-avoiding word $1423$ forces vertex $4$ in BFS order to have parent position $3$, while $1243$ cannot use that parent position. Thus no single unary-binary shape realizes all words in $\UB_4^{(2)}(231)$.

\section{Binary enumeration and sharpness}\label{sec:binary}

\subsection{Sequence identities}
\begin{corollary}\label{cor:oeis}
For $n\ge1$ and $r\ge1$,
\begin{align*}
\seq{A245899}(n)&=\seq{A246747}(n),\\
\seq{A245902}(r)&=\seq{A245899}(2r-1),\\
\seq{A245901}(r)&=\seq{A245898}(2r-1).
\end{align*}
\end{corollary}
\begin{proof}
The first identity follows from Theorem~\ref{thm:A} with $k=2$, since a word determines its labeling on the complete shape. The second is Corollary~\ref{cor:full312}, and the third is Theorem~\ref{thm:B}. These are identities of word sets before taking cardinalities.
\end{proof}

Levin et al.\ establish the relevant heap recurrence and the equinumerosity of $231$- and $312$-avoiding heaps in Theorems~5 and~6 of \cite{LPRS2016}. Transporting that recurrence across Corollary~\ref{cor:oeis} gives
\begin{equation}
a(n)=\sum_{i=0}^{\lfloor(n-1)/2\rfloor} C_i\,a(n-i-1),\qquad n\ge1,\quad a(0)=1,\label{eq:recurrence}
\end{equation}
where $C_i=\frac1{i+1}\binom{2i}{i}$ and $a(n)=\seq{A245899}(n)$ for positive $n$. The value $a(0)=1$ is the auxiliary empty-heap convention; the word-counting OEIS entry begins at $n=1$.

\subsection{Two different obstructions for 321}
\begin{proposition}\label{prop:heapfailure}
The analogue of Theorem~\ref{thm:A} for binary trees and pattern $321$ fails first at length $4$.
\end{proposition}
\begin{proof}
The word $1423$ avoids $321$ and is realized by the edges $1\to4$, $1\to2$, and $2\to3$, with root children ordered $4,2$. The heap edge at position $4$ would require $4<3$. At lengths at most $3$, every nonroot heap parent is the root, whose label is $1$, so no counterexample exists.
\end{proof}

\begin{proposition}\label{prop:parityfailure}
The $321$ full binary collapse first fails at odd length $11$. At that length,
\[
|\UB_{11}^{(2)}(321)|=8095,\qquad |\FB_{11}^{(2)}(321)|=8048.
\]
The lexicographically least word in their difference is
\[
(1,4,2,5,3,7,8,9,10,6,11).
\]
\end{proposition}
\begin{proof}
Exhaustive enumeration, described in Section~\ref{sec:verification}, gives equality of the sets at lengths $1,3,5,7,9$, the two counts at length $11$, and their complete $47$-word difference. It also determines the least word displayed above.

For this particular word, realizability and nonrealizability admit direct certificates. A unary-binary realization has BFS outdegrees
\[
(2,1,2,2,2,0,0,0,0,1,0).
\]
The resulting edges all increase. The word avoids $321$: within its first five positions there is no decreasing triple, the values $7,8,9,10$ increase, the later $6$ exceeds every value in those first five positions, and the terminal $11$ cannot complete a decreasing triple.

In a full binary realization the first levels must be $(1)$, $(4,2)$, and $(5,3)$. Indeed, $4$ cannot take the pair $(5,3)$, so $2$ must take it. The next level has size two or four. Size four gives $(7,8,9,10)$, below which the next value $6$ has no possible parent. Size two gives $(7,8)$; a following level of size four would contain $6$, again impossible. The following level must therefore be $(9,10)$, after which $6$ still has no possible parent. No full binary realization exists.
\end{proof}

\subsection{Additional sequence entries}

Table~\ref{tab:terms} extends the retained baseline of eight terms in each unary-binary entry and five terms in each full binary entry. These $39$ baseline entries are reproduced by both enumeration programs. The table contains $28$ additional entries: $27$ obtained by direct enumeration and $\seq{A245898}(15)$ obtained from the full binary count using Corollary~\ref{cor:oeis}. Here ``additional'' refers to entries in these six sequences relative to that baseline; the $312$ values also occur in the previously known heap sequence.

\begin{table}[ht]
\centering\small
\begin{tabular}{@{}rrrr@{\hskip 1.5em}rrrr@{}}
\toprule
\multicolumn{4}{c}{Unary-binary, length $n$}&\multicolumn{4}{c}{Full binary, length $\ell$}\\
\cmidrule(r{1em}){1-4}\cmidrule{5-8}
$n$&\seq{A245898}&\seq{A245899}&\seq{A245900}&$\ell$&\seq{A245901}&\seq{A245902}&\seq{A245903}\\
&$(231)$&$(312)$&$(321)$&&$(231)$&$(312)$&$(321)$\\
\midrule
9&667&222&753&11&6788&1601&8048\\
10&2098&544&2439&13&74969&12416&92716\\
11&6788&1601&8095&15&877865&105769&1125054\\
12&22349&4095&27350&&&&\\
13&74969&12416&93885&&&&\\
14&254848&33785&326396&&&&\\
15&$877865^{\ast}$&&&&&&\\
\bottomrule
\end{tabular}
\caption{Additional entries relative to the recorded baseline. The starred entry follows from the proved $231$ identity. Unfilled cells are outside the displayed computational range. Directly enumerated terms through length $11$ have two enumeration checks; larger direct terms were recomputed with the prefix-search solver. The $312$ terms also agree with recurrence~\eqref{eq:recurrence}.}\label{tab:terms}
\end{table}

\section{Exponential growth}\label{sec:growth}

\begin{corollary}\label{cor:growth}
For every fixed $k\ge2$ and $\tau\in\{231,312,321\}$,
\[
\lim_{n\to\infty}|\UB_n^{(k)}(\tau)|^{1/n}=4,
\qquad
\lim_{\substack{n\to\infty\\n\equiv1\, (\mathrm{mod}\,k)}}
|\FB_n^{(k)}(\tau)|^{1/n}=4.
\]
\end{corollary}
\begin{proof}
Each of the three patterns is sum indecomposable, meaning it cannot be written as a direct sum of two shorter nonempty permutations. Defant's Theorem~2.1 \cite{Defant2019} therefore gives
$\lim_n|\HP_n^{(k)}(\tau)|^{1/n}=\lim_n|\Av_n(\tau)|^{1/n}=4$,
since the latter counts are Catalan numbers. The inclusions
\[
\HP_n^{(k)}(\tau)\subseteq\UB_n^{(k)}(\tau)\subseteq\Av_n(\tau)
\]
prove the first limit by squeezing. At admissible full-tree sizes the complete shape is full, so the same argument with $\FB$ proves the second.
\end{proof}

These limits determine exponential growth per vertex. In the full binary OEIS indexing, where entry $r$ corresponds to $2r-1$ vertices, the corresponding $r$-th root limit is $16$. Exponential growth does not by itself establish convergence of consecutive ratios or a leading asymptotic term. In particular, the early ratios of \seq{A245898} below $4$ do not identify a smaller limiting growth constant.

\section{Computational and formal verification}\label{sec:verification}

\subsection{Enumeration and certificates}
The reference program recursively assigns zero, one, or two increasing children to each current parent, or zero or two in full binary mode. It collects complete BFS words in a set and filters them by a cubic triple test. This counts each word once even when several trees realize it. The second program enumerates permutation prefixes, discarding those that contain the forbidden pattern or have no feasible level-matching state. Its state includes level boundaries and a nondeterministic parent-child matching automaton. Agreement between the programs checks two different realizability implementations; it does not assume probabilistic independence of their possible errors.

Both programs reproduce all $39$ baseline entries and all $12$ directly computed extensions through length $11$. The complete retained table is recomputed by the second program through unary-binary length $14$ and full binary length $15$. Equation~\eqref{eq:recurrence} checks every displayed $312$ value and generates the longer b-files: indices $1$ through $1000$ for \seq{A245899} and $1$ through $500$ for \seq{A245902}. Every b-file is checked for its index range and values.

The certificate verifier compares the retained list of $47$ words with the exact difference of the reference word sets. It also checks each word using a separate dynamic program over BFS degree sequences, which uses neither level boundaries nor the matching automaton. The same program tests the general-$k$ statements on every permutation through length $7$ for $k=2,3,4$, and reconstructs the BFS traversal produced by the full-degree construction. These are finite computational checks of the implementation and examples; the unrestricted statements are established by the proofs.

The original binary test suite is retained. It checks Theorem~\ref{thm:A} through length $11$, Theorem~\ref{thm:B} at lengths $9$ and $11$, and three auxiliary matching properties. Its $127127$ promotion-cascade instances range over unary-binary-realizable permutations of sizes $8$ and $9$ without a pattern-avoidance filter. The unconditional promotion lemma is proved separately and is not inferred from this finite sample.

\subsection{Lean statements and scope}
The Lean~4 development uses the pinned toolchain and no external library dependencies \cite{Lean4}. The binary results \texttt{theoremA\_trees} and \texttt{theoremB\_trees} quantify over an inductive type of ordered labeled trees. The definitions of increasing labels, unary-binary and full binary degrees, and breadth-first traversal are connected to the original level-matching representation by proved equivalences.

The general-$k$ development verifies the parent-sequence arguments: the capacity bound, the complete-heap inequalities for $312$, and the block construction with increasing labels for $231$. Proposition~\ref{prop:parents} supplies the general tree interpretation in the mathematical proof; the separate binary development supplies the kernel-checked tree interpretation for the OEIS applications. The formalization records these scopes explicitly. Kernel verification establishes the declared mathematical statements; the enumeration counts and the literature attributions are checked by their respective computational and source evidence.

\subsection{Reproduction and availability}
The public repository contains the LaTeX source, Lean sources, both enumerators, the parent-sequence verifier, raw sequence tables, b-files, and certificates. The release includes the PDF, a complete source and evidence archive, a machine-readable manifest, checksums, and execution logs. The following commands reproduce the checks:
\begin{verbatim}
make verify-python   # binary lemmas, words, certificates, b-files
make verify-lean     # kernel checks and axiom audit
make recompute       # all retained direct sequence terms
make paper           # rebuild this PDF
\end{verbatim}
The repository and versioned release are available at
\url{https://github.com/vicotrbb/bfs-avoidance-frontier}.

\section{Further questions}\label{sec:questions}

The structural identities leave several enumeration questions. An exact recurrence or generating function for $\UB_n^{(2)}(231)$ would complement the known heap recurrence for $312$. Corollary~\ref{cor:growth} fixes the exponential rate, so the remaining asymptotic questions concern subexponential factors, possible periodic effects, and convergence of ratios.

For $321$, the difference
$|\UB_{2r-1}^{(2)}(321)|-|\FB_{2r-1}^{(2)}(321)|$
vanishes for $r\le5$ and equals $47$ at $r=6$. A structural characterization of these missing words could replace exhaustive nonrealizability tests with a description of the obstruction. More generally, one can ask which pattern classes admit degree normalization by parent blocks, and which restrictions on allowed outdegrees are preserved by the promotion lemma.

The algorithms also suggest a distinction worth studying: a word may admit a rapidly checkable canonical realization even when direct enumeration over shapes is expensive. For $312$, Theorem~\ref{thm:A} supplies such a recognition criterion. For $231$, Theorem~\ref{thm:B} transforms a supplied realization. A structural recognition criterion simpler than the general degree-sequence dynamic program would complement this construction.

\section{Conclusion}

Breadth-first parent sequences expose two mechanisms behind pattern-avoiding tree realizability. For $312$, the capacity lower bound on parent positions combines with one forbidden triple to force the complete heap inequalities. For $231$, forward compatibility allows consecutive blocks of $k$ children to share a parent, giving a full $k$-ary realization whenever the size permits one. These arguments prove word-set equalities for every branching bound and yield the three binary OEIS identities as specializations.

For binary trees, $321$ heap collapse fails at four vertices and full-degree collapse at eleven. The formal proofs, sequence tables, and complete counterexample set make these distinctions reproducible. Exponential growth is $4$; finer enumeration for $231$ and a structural classification of $321$ obstructions remain open.

\appendix
\section{A pattern-free promotion principle}\label{sec:cascade}

A level matching assigns a consecutive group of children to each parent in order, with at most $k$ children per parent and increasing labels along each edge. A suffix realization below a previous level $P$ is a sequence of nonempty levels with valid adjacent matchings. Define its total depth as $\delta(L_1,\ldots,L_m)=\sum_{t=1}^m t|L_t|$.

\begin{lemma}[Promotion]\label{lem:promotion}
Suppose a nonempty suffix $x u$ is realizable below $P$ with maximum outdegree $k$. Then $u$ is realizable below the extended previous level $Px$, with strictly smaller total depth. No pattern-avoidance or distinctness hypothesis is needed.
\end{lemma}
\begin{proof}
Let $J_0=\{x\}$ and let $J_t$ be the children of $J_{t-1}$ in the old level $L_{t+1}$. Each $J_t$ is an initial segment, because the matchings preserve parent order. Promote $x$ and all its descendants by one level, replacing
\[
L_t\quad\hbox{with}\quad L'_t=(L_t\setminus J_{t-1})J_t.
\]
The new levels remain contiguous blocks of the word. Every remaining edge is an old edge: nonpromoted children stay with nonpromoted parents, and promoted children stay with promoted parents. At the top, the old children of $x$ are appended to the children remaining below $P$. The edge formerly entering $x$ is deleted. Outdegrees therefore remain unchanged except for the parent that loses $x$, whose degree decreases by one. All degree upper bounds and label inequalities persist.

If a new level is empty, then its old level consisted entirely of promoted vertices and their child set is empty. All deeper levels are empty as well, so trailing empty levels are discarded. Every promoted vertex moves up one level, no vertex moves down, and $x$ leaves the suffix. Total depth decreases strictly.
\end{proof}

The lemma works for any fixed outdegree upper bound. It does not by itself preserve the full-degree condition, since the parent of $x$ loses one child. The binary case, including strict depth decrease, is formalized as \texttt{chain\_plus}. The main theorems use the direct parent-sequence constructions instead.

\section*{Acknowledgments and assistance}
The author used Claude (Anthropic) and Codex (OpenAI) for assistance with mathematical exploration, manuscript revision, programming, and Lean formalization. The repository records the mathematical statements and executable verification artifacts. The author is responsible for the definitions, arguments, attribution, and final presentation. The OEIS records and the work of Levin, Pudwell, Riehl, Sandberg, and Defant provided the enumerative questions and the prior results used here.

\clearpage
\begin{thebibliography}{10}
\bibitem{AERRSV2016} J.~Anderson, E.~Egge, M.~Riehl, L.~Ryan, R.~Steinke, and Y.~Vaughan, Pattern avoiding linear extensions of rectangular posets, arXiv:1605.06825, 2016.
\bibitem{BFMR2017} J.~Bettinelli, \'E.~Fusy, C.~Mailler, and L.~Randazzo, A bijective study of basketball walks, \emph{S\'eminaire Lotharingien de Combinatoire} \textbf{77} (2017), Article B77a. \url{https://arxiv.org/abs/1611.01478}.
\bibitem{Bona2012} M.~B\'ona, \emph{Combinatorics of Permutations}, second edition, Chapman and Hall/CRC, 2012.
\bibitem{Defant2019} C.~Defant, Proofs of conjectures about pattern-avoiding linear extensions, \emph{Discrete Mathematics and Theoretical Computer Science} \textbf{21}(4) (2019), \#16. \url{https://doi.org/10.23638/DMTCS-21-4-16}.
\bibitem{Kitaev2011} S.~Kitaev, \emph{Patterns in Permutations and Words}, Springer, 2011.
\bibitem{Knuth1997} D.~E. Knuth, \emph{The Art of Computer Programming, Volume 1: Fundamental Algorithms}, third edition, Addison-Wesley, 1997.
\bibitem{LPRS2016} D.~Levin, L.~K. Pudwell, M.~Riehl, and A.~Sandberg, Pattern avoidance in $k$-ary heaps, \emph{Australasian Journal of Combinatorics} \textbf{64}(1) (2016), 120--139. \url{https://ajc.maths.uq.edu.au/pdf/64/ajc_v64_p120.pdf}.
\bibitem{Lean4} L.~de~Moura and S.~Ullrich, The Lean 4 theorem prover and programming language, in \emph{Automated Deduction, CADE 28}, Lecture Notes in Computer Science \textbf{12699}, Springer, 2021, 625--635.
\bibitem{OEIS} OEIS Foundation Inc., \emph{The On-Line Encyclopedia of Integer Sequences}, \url{https://oeis.org}. Entries \seq{A245898}, \seq{A245899}, \seq{A245900}, \seq{A245901}, \seq{A245902}, \seq{A245903}, and \seq{A246747}; accessed September 2026.
\end{thebibliography}
\end{document}
